\SourceLeaf{038}
\Article{27}
Die Funktion \(\Omega\) lässt sich in unbestimmt vielen Formen darstellen, deren Entwicklung sich lohnen wird. Zunächst erhalten wir durch Quadrieren der Gleichungen I aus Art. 20 und Addieren unmittelbar
\[
\begin{aligned}
\Omega={}&xx\Sigma aa+yy\Sigma bb+zz\Sigma cc+\text{etc.}\\
&+2xy\Sigma ab+2xz\Sigma ac+2yz\Sigma bc+\text{etc.}\\
&+2x\Sigma al+2y\Sigma bl+2z\Sigma cl+\text{etc.}+\Sigma ll
\end{aligned}
\]
Dies ist die \textit{erste} Form.

Multiplizieren wir dieselben Gleichungen jeweils mit \(v,v',v''\) etc. und addieren, so erhalten wir:
\[
\Omega=\xi x+\eta y+\zeta z+\text{etc.}+lv+l'v'+l''v''+\text{etc.}
\]
und daraus, wenn wir für \(v,v',v''\) etc. die im vorigen Artikel angegebenen Ausdrücke einsetzen,
\[
\Omega=\xi x+\eta y+\zeta z+\text{etc.}-A\xi-B\eta-C\zeta-\text{etc.}+M
\]
oder
\[
\Omega=\xi(x-A)+\eta(y-B)+\zeta(z-C)+\text{etc.}+M
\]
Dies ist die \textit{zweite} Form.

Setzen wir in der zweiten Form für \(x-A,y-B,z-C\) etc. die Ausdrücke VII aus Art. 21 ein, so erhalten wir die \textit{dritte} Form:
\[
\begin{aligned}
\Omega={}&[\alpha\alpha]\xi\xi+[\beta\beta]\eta\eta+[\gamma\gamma]\zeta\zeta+\text{etc.}\\
&+2[\alpha\beta]\xi\eta+2[\alpha\gamma]\xi\zeta+2[\beta\gamma]\eta\zeta+\text{etc.}+M
\end{aligned}
\]
Diesen lässt sich die \textit{vierte} Form hinzufügen, die sich aus der dritten Form und den Formeln des vorigen Artikels von selbst ergibt:
\[
\Omega=(v-\lambda)^2+(v'-\lambda')^2+(v''-\lambda'')^2+\text{etc.}+M,\quad\text{oder}
\]
\[
\Omega=M+\Sigma(v-\lambda)^2
\]
Diese Form führt uns die Bedingung für ein Minimum unmittelbar vor Augen.
\Article{28}
Es seien \(e,e',e''\) etc. die Fehler in den Beobachtungen, die \(V=L,\allowbreak V'=L'\quad V''=L''\) etc. ergeben haben; die wahren Werte der Funktionen \(V,V',V''\) etc. seien also jeweils \(L-e,L'-e',L''-e''\) etc. und folglich die wahren Werte von \(v,v',v''\) etc. jeweils \(-e\BareRoot{p},\allowbreak-e'\BareRoot{p'},\allowbreak-e''\BareRoot{p''}\) etc. Der wahre Wert von \(x\) ist daher
\SourceLeaf{039}
\[
=A-\alpha e\BareRoot{p}-\alpha'e'\BareRoot{p'}-\alpha''e''\BareRoot{p''}-\text{etc.}
\]
oder für den bei der besten Bestimmung des Wertes von \(x\) begangenen Fehler, den wir zweckmäßig mit \(Ex\) bezeichnen, ergibt sich
\[
=\alpha e\BareRoot{p}+\alpha'e'\BareRoot{p'}+\alpha''e''\BareRoot{p''}+\text{etc.}
\]
Ebenso ist der bei der besten Bestimmung des Wertes von \(y\) begangene Fehler, den wir mit \(Ey\) bezeichnen werden,
\[
=\beta e\BareRoot{p}+\beta'e'\BareRoot{p'}+\beta''e''\BareRoot{p''}+\text{etc.}
\]
Für den mittleren Wert des Quadrats \((Ex)^2\) ergibt sich
\[
=mmp(\alpha\alpha+\alpha'\alpha'+\alpha''\alpha''+\text{etc.})=mmp[\alpha\alpha]
\]
und für den mittleren Wert des Quadrats \((Ey)^2\) ergibt sich ebenso \(=mmp[\beta\beta]\) etc., wie wir bereits oben gezeigt haben. Nun können wir auch den mittleren Wert des \textit{Produkts} \(Ex.Ey\) angeben; dafür ergibt sich
\[
=mmp(\alpha\beta+\alpha'\beta'+\alpha''\beta''+\text{etc.})=mmp[\alpha\beta]
\]
Dies lässt sich auch in folgender übersichtlicher Form ausdrücken. Die mittleren Werte der Quadrate \((Ex)^2,(Ey)^2\) etc. sind jeweils gleich den Produkten aus \(\frac12mmp\) und den partiellen Ableitungen zweiter Ordnung
\[
\frac{dd\Omega}{d\xi^2},\quad\frac{dd\Omega}{d\eta^2}\quad\text{etc.}
\]
und der mittlere Wert eines Produkts wie \(Ex.Ey\) ist gleich dem Produkt aus \(\frac12mmp\) und der Ableitung \(\frac{dd\Omega}{d\xi.d\eta}\): sofern dabei \(\Omega\) als Funktion der Unbestimmten \(\xi,\eta,\zeta\) etc. betrachtet wird.
\clearpage
